%!TEX root = ../csuthesis_main.tex
% 设置中文摘要
\keywordscn{双目视觉；三维人体姿态估计；轻量化网络；混合精度量化；移动端部署}
\categorycn{TP391}
\itemcountcn{图 \totalfigures\ 幅，表 \totaltables\ 个，参考文献 \total{citnum}\ 篇}

\begin{abstractcn}\setlength{\baselineskip}{20pt}%\renewcommand{\baselinestretch}{1.0}
面向移动机器人、智能安防和沉浸式交互等场景，多人三维人体姿态估计需要同时满足空间定位精度和端侧实时性要求。单目方法易受深度歧义影响，传统双目级联方法又普遍存在对称双分支冗余、稠密深度估计开销大以及移动端部署困难等问题。因此，如何在保持双目三维估计精度的同时实现轻量化建模与高效部署，成为资源受限设备上的关键研究问题。

围绕上述问题，构建了面向移动端的轻量化双目三维人体姿态估计方法。模型层面，设计“左重右轻”的非对称双目骨干网络，以降低双目对称建模带来的冗余计算；提出极线特征增强模块，在一维视差窗口内完成跨视图特征交互，避免显式构建高维代价体；进一步设计可微加权三角化姿态求解器，将软argmax、置信度加权与双目三角化统一到端到端训练框架中，并通过跨分支一致性约束提升轻量分支的几何表征能力。部署层面，面向 Snapdragon 8 Gen3 CPU 平台开展硬件感知优化，通过延迟拆解、Roofline 分析与量化敏感性筛查确定瓶颈层和高风险层，进而设计分层混合精度方案，并采用基于直方图校准的后训练量化方法完成 MNN 模型部署。

在 RT-Pose 数据集上的实验结果表明，所提方法在平均 MPJPE 上达到 158.6\,mm，优于对比的单目基线与基于稠密深度先验的级联范式，同时模型参数量为 11.70\,M，FLOPs 为 13.93\,G。真机部署实验表明，混合精度量化方案在 Xiaomi 14 上可将 FP16 基线的单人场景推理时延由 116.0\,ms 降低至 72.3\,ms，实现 1.60$\times$ 加速，并将模型文件大小由 23.4\,MB 压缩至 15.2\,MB，且仅带来 2.2\,mm 的 MPJPE 增量。多人场景测试进一步表明，量化后的整图推理链路仍保持稳定的时延增长特性。

研究结果表明，面向移动端双目三维人体姿态估计任务，采用非对称结构设计、几何约束驱动的跨视图特征增强以及硬件感知混合精度部署，可以在较好保持三维估计精度的同时显著提升端侧推理效率。相关研究可为资源受限设备上的双目三维感知系统设计与部署提供参考。
\end{abstractcn}
